\chapter{Data and Experimental Protocol}\label{ch:data}

\section{Pretraining corpus (MAESTRO)}
The music model is pretrained on the MAESTRO corpus alone \cite{maestro}.
MAESTRO consists of competition piano performances captured on Disklavier
instruments as finely aligned audio and MIDI, of which only the MIDI is used
here. The music model is a decoder-only Transformer trained from scratch over
note-structured MIDI tokens by next-token prediction. The reference model used
throughout this thesis reaches a validation perplexity of $10.85$.

\section{Evaluation corpus (ASAP) and target extraction}
ASAP \cite{asap} is the only corpus that pairs written scores with
performances, and it supplies the piano targets $y = [\tau, \log r, v]$
through the score-to-performance beat warp $T(\cdot)$ (Eq.~\ref{eq:variables}).
High-quality piano MIDI yields these targets almost exactly, so the
observation noise in the piano setting is negligible ($\Sigma_e \approx 0$).
The inference task is therefore not denoising but \emph{structured
prediction}: the score graph prior is exercised by imputing masked notes from
the rest of the performance.

\section{Strings-and-winds corpus (URMP)}\label{sec:urmp}
The strings-and-winds setting uses \textbf{URMP}, a corpus of 44
multi-instrument chamber pieces assembled from separately recorded monophonic
tracks. Each track comes with its audio, frame-level ground-truth pitch at a
10\,ms hop, note-level onset, pitch and duration annotations, and the score,
and the 13 instruments span strings, woodwind and brass. Three preparatory
measurements govern how the corpus is used, and they are stated here because
each one changes how a later result must be read.

The first measurement calibrates the pitch tracker against the provided
ground-truth pitch before any model is fit. The tracker's median error is
2--5 cents per instrument, an order of magnitude below typical vibrato
extents, and its voicing recall is at least $0.95$. Its per-frame confidence
predicts its own errors: gross octave errors, called slips below, fall from
$9.6\%$ to $0.4\%$ across confidence quintiles. These measurements fix the
estimator chain of Chapter~\ref{ch:phase2}. The estimator's own variances are
used as the observation-noise row, and the lowest confidence quintile is
discarded.

The second measurement concerns the split between development and held-out
pieces, which is fixed at the level of the \emph{composition} and constructed
deterministically without reference to the data. This granularity is forced
rather than merely prudent, because arrangements of one composition share
identical track recordings, so any finer split would place the same audio on
both sides. The composition of the held-out portion and the rules under which
it was evaluated are given in \S\ref{sec:hygiene}.

The third measurement is an absence: no clean intonation ground truth exists,
and this changes what a result means. The targets are estimator outputs, so
held-out recovery measures the model's ability to recover \emph{the
estimator's} targets at masked notes, not its agreement with a gold standard.
The provided pitch curves supply an independent cross-check. The same
estimator run on the ground-truth pitch curves gives targets that are
independent of the tracker, though not of the estimator, and these are called
the reference targets, or the reference centre, in Chapter~\ref{ch:phase2}.
The result available in the strings-and-winds setting is therefore weaker in
kind than the piano result, and it is not reported as though it were the
same. Figure~\ref{fig:vibrato-compare} shows the three layers, truth,
estimate and prediction, on two notes.

\figph{phase2_vibrato_compare_dev.png}{1.0}{Truth, estimate and prediction
at the note level on a development track, with the same fit and mask as
Fig.~\ref{fig:phase2-real}. Top: a held-out note, whose estimate the model
never saw. The estimate of eq.~\eqref{eq:cents-curve} and \eqref{eq:vibrato}
tracks the truth-derived parameters closely (extent $12.7$ against $12.8$
cents). The truth lies inside the model's $90\%$ bands for centre, extent and
rate, while the delay prediction misses (truth $0$\,ms against $136 \pm 48$),
one illustration of why no hypothesis is stated for the delay channel.
Bottom: a visible note whose vibrato the estimator could not identify because
few confidence-kept frames remained. The model's predictions place the truth
within roughly one standard deviation on all three vibrato parameters. The
centre row shows the scope of the setting: on this sparse note the estimator
sits $8$ cents from the truth and the model recovers the estimator's value,
since recovery is agreement with the estimator. Model curves are drawn from
predicted parameters, and phase is not a target.\label{fig:vibrato-compare}}

\section{Contamination control and the strict protocol}\label{sec:contam}
ASAP reuses many MAESTRO recordings, so every ASAP performance whose MAESTRO
twin was in pretraining is removed ($1036 \to 653$ aligned performances), and
the model is tested only on performances the music model has not seen. A
\emph{strict} protocol additionally blanks every hidden note's velocity from
the music model's input, so that no hidden target reaches the model through
the token stream. The cost of the strict protocol is negligible on error,
$+0.0003$ for the music-model system and $+0.0000$ for the feature and
music-model stack, and small on calibration, approximately $+0.014$ in NLL.
The strict numbers are the ones quoted throughout the thesis.

\section{Development and held-out evaluation}\label{sec:hygiene}
Two disjoint sets of pieces play two different roles in every setting of this
thesis. The development set is where every modelling choice was made: the
noise floor, the correction of the music-model read-out, the feature stack,
the kernel and edge families, and the proposed model itself were all selected
on the \nPieces{} development pieces of the piano setting and on the
development side of the composition-level split in the strings-and-winds
setting. Each paired comparison on a development set is valid on its own
terms, but a set that is reused for selection favours whatever survived the
selection, so every number obtained on a development set is labelled as a
development number wherever it appears. In the usual machine-learning
vocabulary the two sets are the validation and test splits \cite{nosek2018}.
They are called the development set and the held-out test set here to make
explicit that the second was consulted only after the modelling choices had
been settled.

Each held-out test set was evaluated once, against hypotheses fixed in
advance. Before the test pieces were examined, the systems to be compared,
the primary and secondary measures, and the decision rule stating which
measure decides each hypothesis and in which direction were written down and
committed to the project's version history. The evaluation was then run a
single time, and every result of that run is reported, including the results
that did not support the hypotheses. The piano held-out test set comprises
\nConfPieces{} further contamination-filtered pieces from the same corpus
shuffle, verified untouched by any modelling decision and evaluated with mask
seeds disjoint from the development seeds. The strings-and-winds held-out
test set comprises 13 of the 44 URMP pieces, chosen at the composition level
(\S\ref{sec:urmp}) and evaluated in Chapter~\ref{ch:phase2}. The inflation
that this procedure protects against is visible in the results themselves:
every system's calibration is worse on the held-out pieces than on the
development set, and one development-set advantage, the pooled NLL, did not
reappear on the piano test set.

The strings-and-winds test set was used a second time, and that second use is
disclosed wherever its results appear. The first evaluation tested the main
hypotheses of the strings-and-winds setting. Two follow-up hypotheses, a
correlated noise model for the timing channel and a tonal pitch metric, were
then developed on the development set alone, with their measures and decision
rules again fixed in advance, and were evaluated once more on the same test
pieces (\S\ref{sec:phase2-reuse}). The consequence for interpretation is that
the second evaluation is a weaker check than the first. The follow-up
hypotheses were formed without access to the test pieces, so the second
evaluation can show that a pre-specified effect appears on pieces that played
no part in its development, but the pieces were no longer untouched, since
their first-round results were known when the follow-up hypotheses were
formed. The two follow-up results are read with that qualification, and any
further use of the same pieces would have to be reported as exploratory.

Every comparison between two systems in every results table follows the same
conventions. The contrast is paired at the level of the piece: the two
systems see identical pieces, identical masks and identical seeds, so their
difference is computed per piece and only then aggregated. The interval
reported with each paired difference comes from a percentile bootstrap over
pieces, in which the per-piece differences are resampled with replacement,
$B = 2000$ times, and the $2.5$th and $97.5$th percentiles of the resampled
mean form the interval. A difference is called significant when that interval
excludes zero, and significance is always stated in those words rather than
by a mark in a table. Pairing is essential here, because piece difficulty
varies far more than the gap between systems, so an unpaired comparison would
obscure every effect reported in this thesis.

\section{Sensitivity to the masking level}\label{sec:masksweep}
The operating point used in the results hides \maskRate{} of the notes, a
level inherited from earlier work rather than derived from the data. A
development-set sweep carried out after the held-out evaluation, from 50\%
hidden down to the leave-one-out limit (Appendix~\ref{app:masksweep}), finds
the ordering of the models stable at every level. The score graph's paired
contribution has an interval that excludes zero from 50\% down to 20\% and
again at leave-one-out, and coverage is insensitive to the level throughout.
The operating point is therefore arbitrary, but the conclusions drawn at it
do not depend on it.

\section{Metrics for recovery and calibration}\label{sec:metrics}
Predictions on held-out notes are scored on both recovery and calibration.
With held-out values $y_i$, predictive means $m_i$, and floored predictive
standard deviations $\sigma_i^{\mathrm{pred}} = \sqrt{(\Sy)_{ii} + \sigma_e^2}$
over $n$ held-out notes, the four measures are
\begin{align}
  \text{RMSE} &= \Bigl( \tfrac1n \textstyle\sum_i (y_i - m_i)^2 \Bigr)^{1/2},
  &
  \text{NLL} &= \tfrac1n \textstyle\sum_i \Bigl[ \tfrac12 \log\bigl(2\pi (\sigma_i^{\mathrm{pred}})^2\bigr)
        + \tfrac{(y_i - m_i)^2}{2 (\sigma_i^{\mathrm{pred}})^2} \Bigr],
  \\[3pt]
  \text{cov}@\gamma &= \tfrac1n \textstyle\sum_i \mathbf{1}\bigl[\, |y_i - m_i| \le z_{\gamma}\, \sigma_i^{\mathrm{pred}} \bigr],
  &
  \text{cal-err} &= \sup_u \bigl|\, \hat F_{\mathrm{PIT}}(u) - u \,\bigr|,
\end{align}
where $z_\gamma$ is the two-sided standard-normal quantile for the level
$\gamma$, which is 0.9 throughout, and $\hat F_{\mathrm{PIT}}$ is the
empirical CDF of the probability integral transform
$u_i = \Phi_{\N}\bigl((y_i - m_i)/\sigma_i^{\mathrm{pred}}\bigr)$, which is
Uniform$(0,1)$ if and only if the predictive distribution is calibrated. Here
$\Phi_{\N}$ is the standard normal CDF, distinct from the waveform synthesis
operator $\Phi(z)$. The pre-specified criterion for a structured prior
requires improvement on \emph{both} recovery (RMSE) and calibration (NLL,
coverage and cal-err) against the independent, temporal-only, ridge and
music-model-mean-only baselines, on held-out, contamination-filtered data. A
method that lowers error but worsens calibration, or the reverse, does not
meet the criterion. Every pooled table additionally reports the median beside
the mean, so that a single collapsed fit cannot hide in the average, and the
safeguard of \S\ref{sec:eb} addresses the cause of such collapses.

% =====================================================================
